\documentclass[12pt]{article}
%packages
\usepackage[top=1in, bottom=1in, left=.75in, right=.75in]{geometry}
\usepackage{amsmath, enumerate, amssymb}
\usepackage{fancyhdr}
\usepackage{graphicx, xcolor, setspace, array}
\usepackage{txfonts}
\usepackage{multicol,coordsys,pgfplots}
\usepackage[scaled=0.86]{helvet}
\usepackage{anyfontsize}
\usepackage{tikz,pgfplots}
\usetikzlibrary{calc,arrows.meta}

% specials
\renewcommand{\emph}[1]{\textsf{\textbf{#1}}}
\def\ra{\rightarrow}
\pgfplotsset{compat = newest}
\newcommand{\blank}[1]{\rule{#1}{0.75pt}}
\pgfplotsset{my style/.append style={axis x line=middle, axis y line=
middle, xlabel={$x$}, ylabel={$y$}}}
\newcommand{\be}{\begin{enumerate}}
\newcommand{\ee}{\end{enumerate}}
\newcommand{\ans}[1][1in]{\rule{#1}{.5pt}}
\newcommand{\ul}[1]{\underline{#1}}
\let\ds\displaystyle
\def\continued{{\emph {Continued....}}}
\def\continuing{{\emph {Problem \arabic{probcount} continued....}}\par\vskip 4pt}

%page setup
\parindent 0pt
\parskip 4pt
\pagestyle{fancy}
\fancyfoot[C]{\emph{\thepage}}
\fancyfoot[R]{} %%%%%% <-- Version Info
\fancyhead[L]{\ifnum \value{page} > 1\relax\emph{Math F113X: Exam 1}\fi}
\fancyhead[R]{\ifnum \value{page} > 1\relax\emph{Fall 2026}\fi}
\headheight 15pt
\renewcommand{\headrulewidth}{0pt}
\renewcommand{\footrulewidth}{0pt}

% Problem Set up 
%% REALLY?
\newcounter{probcount}
\newcounter{subprobcount}
\newcommand{\thesubproblem}{\emph{\alph{subprobcount}.}}
\def\problem#1{\setcounter{subprobcount}{0}%
\addtocounter{probcount}{1}{\emph{\arabic{probcount}.\hskip 1em(#1)}}\par}
\def\subproblem#1{\par\hangindent=1em\hangafter=0{%
\addtocounter{subprobcount}{1}\thesubproblem\emph{#1}\hskip 1em}}
\def\probskip{\vskip 10pt}
\def\medprobskip{\vskip 2in}
\def\subprobskip{\vskip 45pt}
\def\bigprobskip{\vskip 4in}
\newenvironment{subproblems}{%
\begin{enumerate}%
\setcounter{enumi}{\value{subprobcount}}%
\renewcommand{\theenumi}{\emph{\alph{enumi}}}}%
{\setcounter{subprobcount}{\value{enumi}}\end{enumerate}}

%%%%%
% Document Begins
\begin{document}

{\emph{\fontsize{26}{28}\selectfont Fall 2026 \hfill
\hfill Math F113X}}

\begin{center}
{\emph{
{\fontsize{32}{36}\selectfont Math F113X Exam 1}}
}
\end{center}

\vfill
\emph{\fontsize{18}{22}\selectfont Name:\underline{\hspace{2in}}\hspace{2cm} Instructor: \underline{\hspace{2in}}}\\


\vfill
{\fontsize{18}{22}\selectfont\emph{Rules:}}

\begin{itemize}
\item Partial credit will be awarded, but you must show your work.

\item You may have a $3'' \times 5'' $ notecard with writing on both sides.

\item Calculators are allowed. 

\item Turn off anything that might go beep during the exam.

\end{itemize}




Good luck!

\vfill
\def\emptybox{\hbox to 2em{\vrule height 16pt depth 8pt width 0pt\hfil}}
\def\tline{\noalign{\hrule}}
\centerline{\vbox{\offinterlineskip
{
\bf\sf\fontsize{18pt}{22pt}\selectfont
\hrule
\halign{
\vrule#&\strut\quad\hfil#\hfil\quad&\vrule#&\quad\hfil#\hfil\quad
&\vrule#&\quad\hfil#\hfil\quad&\vrule#\cr
height 3pt&\omit&&\omit&&\omit&\cr
&Problem&&Possible&&Score&\cr\tline
height 3pt&\omit&&\omit&&\omit&\cr
&1&& 10 &&\emptybox&\cr\tline
&2&& 16	&&\emptybox&\cr\tline
&3&& 20	&&\emptybox&\cr\tline
&4&& 8	&&\emptybox&\cr\tline
&5&& 12	&&\emptybox&\cr\tline
&6&& 15	&&\emptybox&\cr\tline
&7&& 19	&&\emptybox&\cr\tline
%&8&&	&&\emptybox&\cr\tline
%&9&&	&&\emptybox&\cr\tline \tline \tline
&Extra Credit&&(5)&&\emptybox&\cr\tline\tline \tline
&Total&&100&&\emptybox&\cr
}\hrule}}}

\newpage

\problem{10 points} A section of  Math F113X with 9 students holds a vote for their
favorite vegetable, with the choices being asparagus (A), beets (B) and
cabbage (C). The results are as follows:
\begin{center}
  \begin{tabular}{c|ccccccccc}
  	
	& \rotatebox{90}{Alice} & \rotatebox{90}{Bob} & \rotatebox{90}{Cathy} & \rotatebox{90}{Daniel} & \rotatebox{90}{Edith} & \rotatebox{90}{Fahad} & \rotatebox{90}{Greg} &  \rotatebox{90}{Hyacinth} & \rotatebox{90}{Ida}  \\ \hline
    \textbf{1st choice} & A & A & A & A & A  & B & B & B & C \\
    \textbf{2nd choice} & B & B & C & C & C  & A & A & A & A \\
    \textbf{3rd choice} & C & C & B & B & B  & C & C & C & B
  \end{tabular}
\end{center}
    
    
\begin{subproblems}
\item (4 points) Create a preference schedule with this information.
  \vfill

\item (3 points) How many votes are needed for a majority winner? \ans \\
  Provide a brief calculation/explanation for your answer.
  \vspace{3cm}
  
\item (3 points) What is the smallest number of votes that could give a
  candidate a plurality? \ans \\
  Provide a brief calculation/explanation for your answer.
  \vspace{3cm}
  
\end{subproblems}

\newpage

\problem{16 points} Consider the following preference schedule for an
  election with 5 candidates and 430 voters.
  
  \begin{center}
\begin{tabular}{c|ccccccc}
\textbf{Rank} & \textbf{120} & \textbf{90} & \textbf{80} & \textbf{67} & \textbf{8} & \textbf{62} & \textbf{3} \\
\hline
1st & A & B & C & D & D & E & E \\
2nd & B & A & B & B & C & B & D \\
3rd & C & C & A & C & B & D & B \\
4th & D & D & D & A & A & C & C \\
5th & E & E & E & E & E & A & A \\
\end{tabular}
\end{center}
  
  
Answer the following questions. \emph{Show your work}. (This means write down calculations or words so that someone else can understand how you get your answer.)
\begin{subproblems}



\item (4 points) Find the winner using \emph{Plurality voting}. \ans \\
  Provide a brief calculation/explanation for your answer.
  \vspace{1cm}
  
  \item (2 points) How many votes are needed for a \emph{majority}? \ans \\
  Provide a brief calculation/explanation for your answer.
  \vspace{1cm}
  
\item (10 points) Find the winner using \emph{Instant Runoff Voting} (Ranked Choice
  Voting). \ans \\
  Clearly show your work for each round and \emph{state who gets eliminated} in each round.
  \vfill
  
  
  \end{subproblems}
%\continued
  \newpage
%  \continuing
%  Here's the preference schedule again.

\newcommand{\probTwoPref}{
\begin{tabular}{c|cccc}
\textbf{Rank} & \textbf{6} & \textbf{6} & \textbf{3} & \textbf{9} \\
\hline
1st & A & B & A & C \\
2nd & B & D & D & A \\
3rd & C & C & C & B \\
4th & D & A & B & D \\
\end{tabular}
}

%\begin{center}
%  \probTwoPref
%  \end{center}

\problem{20 points} Consider the following preference schedule for an
  election with 4 candidates and 24 voters.
  
  \begin{center}
  \probTwoPref
  \end{center}

\begin{subproblems}

\item (6 points) Find the winner using \emph{Borda Count}. \ans \\
  Clearly show your work.
  \vfill
  \vfill
  
\item (6 points) Fill in the missing information in the table below and determine
  the winner using \emph{Copeland's method}.
\begin{center}
\begin{tabular}{c|c|c|c|c|c}
\textbf{A vs B} & \textbf{A vs C} & \textbf{A vs D} &
\textbf{B vs C} & \textbf{B vs D} & \textbf{C vs D} \\ \hline
\begin{tabular}{c}
A: {\hspace{0.8cm}}\\
B: 
\end{tabular}
&
\begin{tabular}{c}
A: 9\\
C: 15
\end{tabular}
&
\begin{tabular}{c}
A: 18\\
D: 6
\end{tabular}
&
\begin{tabular}{c}
B: 12\\
C: 12
\end{tabular}
&
\begin{tabular}{c}
B: 21\\
D: 3
\end{tabular}
&
\begin{tabular}{c}
C: {\hspace{0.8cm}}\\
D: 
\end{tabular}
\\ \hline
 &  &  &  &  & 
\end{tabular}
\end{center}

 %A vs B winner: \hfill  C vs D winner: \hfill
\hfill \begin{tabular}[t]{l l l}
Tally:& A: &\ans \\[6 pt]
&B: & \ans \\[6 pt]
& C: & \ans \\[6 pt]
& D: & \ans \\[6 pt]
\end{tabular}

Winner under Copeland's method: \ans

\item (4 points) Is there a \emph{Condorcet winner}? If so, who is it? Briefly
  explain/justify your answer.
  \vfill
  
\item (4 points) Did either Borda Count or Copeland's Method violate the \emph{Condorcet criterion} in this
  situation? If so, which one(s)? Provide a brief explanation of your answer.
  \vfill
  
\end{subproblems}

\newpage

\problem{8 points} %Between Plurality, IRV/RCV, Borda Count, and
%Copeland's Method; choose your favorite voting method (there are no
%wrong answers for this part!). 
Choose one of the following methods: \emph{Plurality, IRV/RCV, Borda Count,} and \emph{
Copeland's Method.}

State at least one \emph{strength} your chosen method has
over the other three options. State at least one \emph{weakness} your chosen method
has. 

You should refer to some of the various \emph{fairness conditions} in discussing the weakness of your method, and explain what the fairness condition you are referring to means in the context of your method.  %that is not shared by the other three methods.

\bigskip

Chosen method: \ans


Strength (4 points):
\vspace{3cm}

Weakness (4 points):
\vspace{3cm}

%\newpage

\problem{12 points}
In a certain corporation, shareholders receive 1 vote for each share of stock they hold.  A small corporation has two people who have 30 shares each, two people who have 20 shares each, and one person who has only 10 shares. 

\begin{subproblems}
 \item (3 points) If any decisions require a majority vote, set up a weighted voting system of the form
 $[q: w_{1}, \ldots, w_{n}]  $
 to represent this corporation's shareholder votes.
 \vfill
 
  \item (3 points) Certain decisions, like amending the bylaws, require a 2/3 vote. How would you change your weighted voting system to represent this situation?
  
  \vfill
  
  \item (3 points)  What is the largest value the quota could be for this system? \ans
    \vfill
  
  \item (3 points) What is the smallest value the quota could be for this system? \ans
    \vfill
  
  \end{subproblems}



%\problem{12 points}
%
%Consider the weighted voting system $[q: 10, 6, 3]$
%
%\begin{subproblems}
%\item What is the largest value that $q$ can be? \ans
%
%\item What is the smallest value that $q$ can be? \ans
%
%\item What should $q$ be if the system has a dictator? $q = $\ans
%
%\item If the quota is 12, 
%
%\be[(i)]
%\item Who has veto power? Explain how you know.
%
%
%
%\item Are there any dummies? Explain how you know.
%
%\ee
%
%\end{subproblems}

\newpage

\problem{15 points} 
Consider the weighted voting system [18: 8,5,4,3,3,2]. Determine:

\begin{subproblems}

\item (1 point) The total number of players:  \ans

\item (1 point)The total weight of the system: \ans
 
\item (1 point) The quota: \ans

\item (3 points) Are there any dictators? Explain how you know, with a few words.

\vfill

\item (3 points) The coalition $\{ P_{1}, P_{2}, P_{3}, P_{6}\}$ is a winning coalition. Explain how you know it is a winning coalition using a computation and some words.

\vfill

\item (3 points) In the coalition $\{ P_{1}, P_{2}, P_{3}, P_{6}\}$, is $P_{6}$ critical? Explain why or why not, using a computation and some words.

\vfill



\item (3 points) Does player $P_{2}$ have veto power? Why or why not?

\vfill

\end{subproblems}

%%%%%%%%

\newpage

\problem{19 points} Consider the weighted voting system $ [27: 19, 18, 8, 2]$ and the list of all possible coalitions.

\begin{center}
{\renewcommand{\arraystretch}{2.}
\begin{tabular}{c | c | c || c | c | c}
coalition & weight & winning? & coalition & weight & winning?\\ \hline \hline
$P_1 \ P_2$ & 37 & & 	$P_1\ P_2\ P_3$ 		&  45 & \\ \hline 
$P_1\ P_3$ & 27 & & 	$P_1\ P_2\ P_4$ 		&  39 & \\ \hline 
$P_1\ P_4$ & 21 & & 	$P_1\ P_3\ P_4$ 		& 29 & \\ \hline
$P_2\ P_3$ &  26 & & 	$P_2\ P_3\ P_4$ 		& 28  & \\ \hline
$P_2\ P_4$ & 20 & & 	$P_1\ P_2\ P_3\ P_4$ 	& 47 &  \\ \hline
$P_3\ P_4$ & 10 &  \\
\end{tabular}
}
\end{center}




\begin{subproblems}
    
    \item (2 points) Put a check $\checkmark$ in the column labeled ``winning?'' to identify all winning coalitions, and mark  $\times$ for non-winning coalitions.

    \item (6 points) In each winning coalition, \emph{underline the critical players}.
    
  \item (5 points) Using this information, compute the \emph{Banzhaf Power Index} of each player by filling in the following table.




%{\renewcommand{\arraystretch}{2.5}
\begin{tabular}{c | p{.75in} | c c }
Player & &  power \\ \hline
&&&\\
&&&\\&&&\\&&&\\&&&\\&&&\\&&&\\&&&\\
\end{tabular}
%}
\item (2 points) Is there a \emph{dummy}? Why or why not?
\vfill

  \item (4 points) Players 1 and 2 have almost the same number of votes, and Player 3 has less than half as many votes as Player 2. Does their decision-making power reflect this vote distribution? Why or why not/how can you tell?
  
  \vfill
   %Does the fact that $P_{2}$ has more votes than $P_{3}$ matter in the decisionmaking represented by this weighted voting system? Why or why not?

\end{subproblems}

%%%%%%%%%
\newpage

\fbox{Extra Credit} (5 points)


A professional basketball team has five coaches: a head coach (H) and four assistant coaches (A1, A2, A3, and A4). 

\begin{quote}
\emph{Constraint:} Player personnel decisions require at least three Yes votes, one of which must be from the head coach (H).
\end{quote}

If $h$ is the weight assigned to the head coach $H$, and $a$ is the weight assigned to each of the four assistant coaches (assuming they each get the same number of votes), determine values  \[ [q: h, a, a, a, a] \] that describe the weighted voting system. Explain why your assignment of values matches the constraint given above. 


\vfill 





%which players are dictators, have veto power, and are dummies:
%if the quota is 10?  
%if the quota is 12?  
%if the quota is 16?


\end{document}
